% !TEX root = ../main.tex
% =====================================================================
% Chapter 5 (Evaluation) --- DISSERTATION version (single-column report class).
% Written from the repository's canonical result documents; every number is
% sourced (no numbers from memory):
%   - docs/reviews/project_review_202607.md  §1.4  (canonical number memo)
%   - docs/analysis/benchmark_mmbv_final_analysis.md  (closing analysis)
%   - docs/results/benchmark_mmbv_final_gpt4judge.md  (paper-setting dimension tables)
%   - docs/results/benchmark_mmbv_final_official_agg.md  (qwen-max robustness control)
%   - docs/results/benchmark_mmbv_dense_{16f,32f}*.md  (frame-scaling)
%   - docs/results/benchmark_v2_agqa.md, annotation_audit.json
%   - docs/progress.md  §14.1 / §14.2 / §14.3 / §16.1 / §16.2
%
% Judge-setting discipline (from thesis_outline.md §4):
%   * paper-level and dimension-level citations = gpt-4-turbo re-judge
%     (1.978 / 1.713 / 1.491, gap 0.265; HL 2.29 vs 0.78);
%   * the runs=3 noise-robustness argument = qwen-max as-run scores with ±std
%     (1.984±0.101 / 1.727±0.020 / 1.478±0.025, gap 0.257 > std-sum 0.121);
%   * duration buckets and the attribution ladder are derived from qwen-max
%     raw scores (stated in the text).
% Honesty red-lines: frames/question is 3.6 (runs=3 authority; 3.1 was runs=1);
% the 81/69 zero-explore split is a runs=1 diagnostic figure; TR is a LOSING
% dimension and is disclosed as such; total tokens are ~2.1x the baseline (no
% cost-saving claims); the 150-question subset is NOT comparable to the public
% leaderboard; frame-scaling main protocol is common-147 (moderation-blocked
% items removed from every condition), full-150 in the appendix.
%
% Cross-references into other chapters: ch:method (exists), ch:diagnosis
% (Ch.3, not yet written -> renders "??" until then).
% =====================================================================

\chapter{Evaluation}
\label{ch:eval}

This chapter evaluates the system of \cref{ch:method} on MMBench-Video with a
protocol designed to answer three questions: whether the agent outperforms
direct-VLM baselines under honest accounting (\cref{sec:main-results}), where the
gain comes from (\cref{sec:attribution}), and what the perception budget buys
(\cref{sec:budget}). Analyses by duration and capability dimension
(\cref{sec:analysis}), falsified alternatives (\cref{sec:negative}), a transfer
check (\cref{sec:agqa}), an annotation audit with threats to validity
(\cref{sec:threats}), and two worked traces (\cref{sec:case}) complete the picture.

\section{Experimental Setup}
\label{sec:setup}

\textbf{Benchmark.} We evaluate on a 150-question subset of
MMBench-Video~\cite{fang2024mmbenchvideo} spanning 136 videos, drawn by stratified
sampling (seed 42) with deliberate oversampling of the Temporal Reasoning (TR) and
Hallucination (HL) dimensions. Because the subset re-balances the official
distribution, \emph{our scores are not directly comparable to the public
leaderboard}, which is computed on the full 1{,}998 questions; re-weighting our
per-dimension scores by the official distribution moves every method's total by at
most $0.04$ (agent $+0.017$), so the subset composition does not drive any
conclusion below.

\textbf{Scoring protocol.} We replicate the official
VLMEvalKit~\cite{duan2024vlmevalkit} protocol: an LLM judge scores each answer on a
0--3 scale (judge prompt, scale, temperature $0$, and failure handling replicated
verbatim), and dimension aggregation follows the official multi-label
\texttt{get\_dimension\_rating} (each question counts toward every capability it is
tagged with: 26 leaf capabilities, 9 L2 dimensions, and Perception/Reasoning/Overall
roll-ups). Two judges are used throughout. The as-run judge is
\texttt{qwen-max}; after the campaign, every cached answer was re-judged by
\texttt{gpt-4-turbo}~\cite{openai2023gpt4} --- the judge family used by the official
benchmark --- totalling 1{,}350 judgements. Unless stated otherwise,
\emph{paper-level and dimension-level numbers below use the gpt-4-turbo setting},
with qwen-max retained as a robustness control; the two settings agree to within
$0.015$ on every method's total, with per-question agreement 0.76--0.81. Derived
analyses (duration buckets, the attribution ladder) are computed from the qwen-max
raw scores and labelled as such.

\textbf{Systems.} Four systems are compared, all sharing the same backbone LLM
(\texttt{qwen-plus-latest}~\cite{bai2023qwen}) and VLM
(\texttt{qwen-vl-plus}~\cite{bai2023qwenvl}; cf.\ \cref{sec:impl}):
\begin{itemize}
  \item \textbf{agent\_v2} (ours): the lazy-memory, confidence-driven agent of
  \cref{ch:method}.
  \item \textbf{vlm\_direct@8}: the VLM answers directly from 8 uniformly sampled
  frames.
  \item \textbf{vlm\_transcript@8} (\emph{fair baseline}): identical to
  vlm\_direct@8 \emph{plus} the full ASR transcript in the prompt. Since the agent
  also consumes ASR, this baseline holds the input modalities fixed and isolates
  the contribution of the architecture (\cref{sec:attribution}).
  \item \textbf{v1}: the scene-graph-centric baseline diagnosed in
  \cref{ch:diagnosis} (legacy, single run).
\end{itemize}

\textbf{Runs and accounting.} The three main systems are run three times
independently (runs=3); we report mean $\pm$ standard deviation. All token, call,
and frame counts come from a thread-safe ledger recording the API-reported usage of
every call --- no estimates (\cref{sec:impl}). Beyond accuracy we report
\emph{frames-touched per question}: the number of frames the system actually sent
to the VLM for that question.

\section{Main Results}
\label{sec:main-results}

\Cref{tab:main} shows the headline result: the agent outperforms both baselines
under both judge settings.

\begin{table}[t]
  \centering
  \small
  \caption{Overall results on the MMBench-Video subset (0--3 scale, mean $\pm$ std
  over 3 runs). The gpt-4-turbo column is the paper-level setting; qwen-max is the
  as-run robustness control. v1 is the legacy baseline of \cref{ch:diagnosis}
  (single run, qwen-max).}
  \label{tab:main}
  \begin{tabular}{@{}lccc@{}}
    \toprule
    Method & Overall (gpt-4-turbo) & Overall (qwen-max) & Frames/Q \\
    \midrule
    \textbf{agent\_v2 (ours)} & $\mathbf{1.978 \pm 0.097}$ & $\mathbf{1.984 \pm 0.101}$ & \textbf{3.6} \\
    vlm\_transcript@8 & $1.713 \pm 0.019$ & $1.727 \pm 0.020$ & 8.0 \\
    vlm\_direct@8 & $1.491 \pm 0.019$ & $1.478 \pm 0.025$ & 8.0 \\
    v1 (scene-graph pipeline) & --- & 1.193 & --- \\
    \bottomrule
  \end{tabular}
\end{table}

\textbf{The reversal survives run noise.} Under the qwen-max setting the gap over
the fair baseline is $1.984 - 1.727 = 0.257$, which exceeds the \emph{sum} of the
two standard deviations ($0.101 + 0.020 = 0.121$); under gpt-4-turbo the gap widens
slightly to $0.265$. The agent's larger run-to-run variance ($\pm 0.101$ against
$\pm 0.020$ for the deterministic-sampling baselines) stems from exploration-path
randomness: at low temperature the agent may choose different windows across runs
(a concrete instance appears in \cref{sec:case}). Even so, the worst agent run
exceeds the best fair-baseline run.

\textbf{v1 confirms the diagnosis.} The scene-graph-centric v1 scores 1.193 ---
below even the 8-frame direct baseline --- reproducing the failure analysed in
\cref{ch:diagnosis}: when triplets are the \emph{answer source}, the graph is a
lossy bottleneck rather than an asset. The v2 reversal therefore cannot be
attributed to ``using a scene graph'' per se, but to using it as an index
(\cref{subsec:provenance}).

\section{Attribution: Modality versus Architecture}
\label{sec:attribution}

The fair baseline decomposes the total gain into two steps (qwen-max raw scores;
the architecture step is $0.265$ under gpt-4-turbo):

\begin{table}[h]
  \centering
  \small
  \caption{Attribution ladder (qwen-max, runs=3 means). Each step changes exactly
  one factor.}
  \label{tab:attribution}
  \begin{tabular}{@{}lccl@{}}
    \toprule
    Step & From $\to$ To & $\Delta$ & Factor isolated \\
    \midrule
    ASR modality & $1.478 \to 1.727$ & $+0.249$ & add narration text, same frames \\
    Architecture & $1.727 \to 1.984$ & $+0.257$ & same modalities, agentic memory + orchestration \\
    \bottomrule
  \end{tabular}
\end{table}

The two contributions are almost exactly equal. This rules out the deflationary
reading that the agent ``wins only because it hears the narration'': granting the
baseline the same narration still leaves a $+0.257$ gap that only the architecture
--- lazy evidence memory, provenance pull-back, and confidence-driven exploration
--- can explain. Conversely, the modality step confirms that ASR is a genuinely
load-bearing input for this benchmark, which \cref{sec:analysis} revisits when
interpreting the TR dimension.

\section{Analysis by Duration and Dimension}
\label{sec:analysis}

\subsection{Video duration: the comfort zone ends at 90 seconds}
\label{subsec:duration}

\begin{table}[t]
  \centering
  \small
  \caption{Overall score by video-duration bucket (qwen-max raw scores, runs=3
  means).}
  \label{tab:duration}
  \begin{tabular}{@{}lccc@{}}
    \toprule
    Duration bucket & agent\_v2 & vlm\_transcript@8 & vlm\_direct@8 \\
    \midrule
    $<$90\,s \hfill ($n=52$) & 1.81 & 1.76 & 1.57 \\
    90--180\,s \hfill ($n=50$) & \textbf{2.10} & 1.55 & 1.34 \\
    $>$180\,s \hfill ($n=48$) & \textbf{2.05} & 1.87 & 1.52 \\
    \bottomrule
  \end{tabular}
\end{table}

\Cref{tab:duration} locates \emph{where} the architecture pays. Below 90 seconds
the agent and the fair baseline are statistically indistinguishable (1.81 vs.\
1.76): a short video fits comfortably into 8 uniform frames plus narration, so
guided exploration has little to add. The picture inverts beyond 90 seconds ---
$2.10$ vs.\ $1.55$ at 90--180\,s and $2.05$ vs.\ $1.87$ beyond 180\,s. This is the
expected signature of the design: uniform sampling dilutes as duration grows, while
question-guided window selection does not. We accordingly describe the system's
contribution as \emph{long-video evidence localization and integration}, with
$\approx$90\,s as the boundary of the baselines' comfort zone.

\subsection{Capability dimensions}
\label{subsec:dimensions}

\begin{table}[t]
  \centering
  \small
  \caption{Per-dimension results in the MMBench-Video Table-3 layout
  (gpt-4-turbo judge, official multi-label aggregation, 0--3). \emph{P.\ Mean} /
  \emph{R.\ Mean} are the Perception / Reasoning roll-ups. Full per-dimension
  tables with standard deviations for both judge settings appear in
  \cref{app:tables}.}
  \label{tab:dimensions}
  \footnotesize
  \setlength{\tabcolsep}{3.2pt}
  \begin{tabular}{@{}lcccccccccccc@{}}
    \toprule
    Model & Overall & CP & FP-S & FP-C & HL & \emph{P. Mean} & LR & AR & RR & CSR & TR & \emph{R. Mean} \\
    \midrule
    agent\_v2 & \textbf{1.98} & \textbf{2.18} & \textbf{2.08} & \textbf{1.47} & \textbf{2.29} & \textbf{2.04} & \textbf{2.18} & 2.07 & 2.07 & 2.21 & 1.74 & 1.98 \\
    vlm\_transcript & 1.71 & 1.78 & 1.81 & 0.90 & 0.78 & 1.42 & 1.73 & \textbf{2.50} & \textbf{2.12} & \textbf{2.28} & \textbf{1.87} & \textbf{2.07} \\
    vlm\_direct & 1.49 & 1.53 & 1.64 & 0.77 & 0.64 & 1.28 & 1.39 & 2.21 & 1.81 & 2.26 & 1.52 & 1.79 \\
    \bottomrule
  \end{tabular}
\end{table}

\Cref{tab:dimensions} decomposes the gain. Three observations:

\textbf{The architectural gain is a perception gain.} The agent wins the Perception
roll-up decisively (2.04 vs.\ 1.42) while the Reasoning roll-up stays near parity
(1.98 vs.\ 2.07, a small loss): all three systems share the same backbone LLM, so
the architecture buys \emph{seeing accurately}, not thinking harder. The largest
single-dimension effects are exactly where verbatim, targeted evidence should
matter: cross-instance fine-grained perception (FP-C, 1.47 vs.\ 0.90) and
hallucination resistance below.

\textbf{Hallucination resistance is the standout mechanism.} On the HL dimension
the agent scores 2.29 against 0.78 for the fair baseline
($\approx$2.9$\times$; gpt-4-turbo setting). Under the qwen-max setting the same
comparison reads $2.42 \pm 0.19$ vs.\ $1.04$ vs.\ $0.62$
($\approx$2.3$\times$ / 3.9$\times$) --- the direction is judge-independent. We
attribute this to the grounding rules of \cref{subsec:grounding}: answering ``no''
about unexplored content is procedurally disallowed, so fabricating an answer is
replaced by either exploration or a stated lack of evidence
(\cref{sec:case} shows one such trace).

\textbf{TR is honestly lost --- and the loss is informative.} On Temporal Reasoning
the agent scores 1.74 against the fair baseline's 1.87. Note the baseline itself
jumps from 1.52 (no ASR) to 1.87 (with ASR): most of the benchmark's temporal
signal lives in the narration, i.e.\ the TR gain is \emph{modality-driven}, not
architecture-driven. This is why we frame the contribution as evidence
localization and integration rather than ``improved temporal reasoning'': on
questions whose temporal structure is already spelled out in the transcript, an
extra retrieval layer can only add noise.

\section{Perception Budget and Frame Scaling}
\label{sec:budget}

\subsection{What the agent spends}
\label{subsec:cost}

\begin{table}[t]
  \centering
  \small
  \caption{Per-question cost accounting (runs=3 means; tokens from API-reported
  usage). Prebuild = once-per-video $L_0$ initialisation amortised per question.
  Wall-clock times are from the as-run campaign.}
  \label{tab:cost}
  \footnotesize
  \setlength{\tabcolsep}{4pt}
  \begin{tabular}{@{}lcccccc@{}}
    \toprule
    Method & Tool calls & Frames/Q & Tokens/Q (answer) & + Prebuild/Q & = Total/Q & Time (s) \\
    \midrule
    agent\_v2 & 1.7 & \textbf{3.6} & 10{,}546 & 6{,}108 & 16{,}655 & 21.8 \\
    vlm\_transcript@8 & --- & 8.0 & 7{,}299 & --- & 7{,}299 & 3.4 \\
    vlm\_direct@8 & --- & 8.0 & 6{,}748 & --- & 6{,}748 & 3.7 \\
    \bottomrule
  \end{tabular}
\end{table}

\Cref{tab:cost} states the price of the reversal honestly. The agent touches
\textbf{3.6 frames per question} against the baselines' fixed 8 --- a 55\% frame
reduction --- and averages only 1.7 tool calls, evidence that the instructed budget
of \cref{subsec:budget} is respected in practice. In exchange, total token
consumption is $\approx$2.1$\times$ the fair baseline (16{,}655 vs.\ 7{,}299 per
question) and latency is $\approx$6$\times$ (21.8\,s vs.\ 3.4\,s): the agentic loop
pays in LLM tokens and wall-clock time for what it saves in perception. The value
proposition is therefore \emph{frame efficiency and long-video accuracy}, not cost
reduction.

In a single-run diagnostic breakdown, 81 of 150 questions were answered from free
retrieval alone (zero exploration) while 69 self-escalated to exploration; the
zero-exploration batch scores far higher on average (2.259 vs.\ 1.551), a gap that
\cref{sec:negative} shows to be selection bias rather than evidence that
exploration hurts.

\subsection{Frame scaling: can the baselines just look at more frames?}
\label{subsec:framescaling}

A natural objection is that the baselines are starved: give a dense sampler enough
frames and it should catch up. We test this directly by densifying both baselines
to 16 and 32 uniform frames (single run each, verbose answer mode) and judging
under both settings.

\textbf{Protocol (common-147).} Dense sampling has a hidden failure mode: three
questions (mmbv-0131, 0679, 1185) are killed by the hosting provider's content
moderation with a deterministic HTTP 400, and the error count rises monotonically
with frame count (0 errors at 8f, 1 at 16f, 3 at 32f; 100\% reproducible on
confirmatory reruns, retries ineffective --- the 8-frame baselines and the agent
score normally on all three). To avoid crediting the agent for infrastructure
failures, we remove these three questions from \emph{every} condition, including
the agent and the 8f baselines, yielding the common-147 protocol; the full-150
accounting, which scores moderation kills as 0 and therefore \emph{underestimates}
the dense baselines, is reported in \cref{app:tables}. Because the dense runs are
single-run, the curves use each system's first trial; the agent's three-run means
(qwen 1.968, gpt-4-turbo 1.957) agree with its trial-1 value in ordering.

\begin{figure}[t]
  \centering
  % Figure is a PDF pre-converted from fig/frame_scaling.svg (rsvg-convert -f pdf).
  % Using the PDF avoids invoking Inkscape on every compile: much faster on
  % Overleaf, and it works in any TeX install without shell-escape.
  % SVG fallback: keep \usepackage{svg} and swap this line back to
  %   \includesvg[inkscapelatex=false,width=0.8\linewidth]{fig/frame_scaling}
  \includegraphics[width=0.8\linewidth]{fig/frame_scaling}
  \caption{Frame-scaling on the MMBench-Video subset (common-147, trial-1): dense
  uniform sampling never overtakes guided sampling under either judge.}
  \label{fig:framescaling}
\end{figure}

\begin{table}[t]
  \centering
  \small
  \caption{Frame-scaling results (0--3, common-147, trial-1). The agent column is
  its $\sim$3.6-frame operating point; the last column is the gap remaining at 32
  frames.}
  \label{tab:framescaling}
  \begin{tabular}{@{}llccccc@{}}
    \toprule
    Judge & Method & 8f & 16f & 32f & agent\_v2 ($\sim$3.6f) & 32f gap \\
    \midrule
    qwen-max & vlm\_direct & 1.442 & 1.497 & 1.680 & \textbf{2.109} & $-0.429$ \\
    qwen-max & vlm\_transcript & 1.735 & 1.660 & 1.823 & \textbf{2.109} & $-0.286$ \\
    gpt-4-turbo & vlm\_direct & 1.510 & 1.707 & 1.857 & \textbf{2.088} & $-0.231$ \\
    gpt-4-turbo & vlm\_transcript & 1.694 & 1.871 & 1.932 & \textbf{2.088} & $-0.156$ \\
    \bottomrule
  \end{tabular}
\end{table}

\textbf{Dense sampling does not cross the curve
(\cref{fig:framescaling,tab:framescaling}).} vlm\_direct rises monotonically with
frames ($+0.12$--$0.17$ per doubling) but at 32 frames still trails the agent by
$0.23$--$0.43$ depending on the judge. Extrapolating the log-linear trend,
vlm\_direct would need roughly 81 frames (gpt-4-turbo) to 389 frames (qwen-max) to
match what the agent achieves at 3.6 --- a \textbf{22$\times$--108$\times$ frame
budget}. The transcript baseline comes closest under the more lenient gpt-4-turbo
judge, where its 32f score (1.932) lands inside the agent's three-run noise band
(1.957, gap 0.025, per-question agreement 0.71); we disclose this as a near
crossover, noting that it requires the full narration modality \emph{plus} a
$\sim$9$\times$ frame budget, and that under the stricter qwen-max judge the same
configuration still trails by 0.29. The answer to ``why not just feed the model
more of the video'' is thus empirical, not speculative: dense sampling sits on a
much worse cost--accuracy frontier, and additionally becomes more fragile as
density grows (the moderation failure mode above).

\section{Ablations and Negative Results}
\label{sec:negative}

The remaining gap to the annotation ceiling (\cref{sec:threats}) invites cheap
fixes. We tested the two most plausible ones; both are falsified, which we regard
as an attribution result: it localises the residual gap in model capability rather
than in the orchestration policy.

\textbf{Oracle routing is falsified by selection bias.} The single-run observation
that zero-exploration questions score higher (2.259 vs.\ 1.551,
\cref{subsec:cost}) suggests routing: skip exploration for dimensions that seem
not to need it. We implemented the \emph{oracle} version --- routing on gold
dimension labels (TR, FP-C, LR), an upper bound no learned router can beat. The
total \emph{drops} from 1.933 to 1.733, every routed dimension falls
(TR $-0.37$, FP-C $-0.27$, LR $-0.40$), and a matched-question comparison shows 2
improved against 13 worsened (40 unchanged). The zero-exploration batch scores
higher because \emph{easier questions do not trigger exploration}, not because
exploration hurts; with even the oracle upper bound negative, the
routing/question-classifier direction is closed.

\textbf{Deeper evidence does not lift the failures.} For the failure cluster
attributed to comprehension rather than retrieval, we raised the per-window frame
cap from 6 to 9 and restructured the caption prompt (intent, interaction,
causality, object enumeration). Zero of the four target questions improved (all
remained at score 0; an unaffected control question held its score), so the change
was reverted --- the bottleneck is VLM/LLM capability, not evidence depth.
Two free diagnostics from the same phase: only 1 of 150 answers exhibited
``hesitation'' (no terminal answer), and no transcript exceeded the context budget,
ruling out two further cheap explanations.

A qualitative taxonomy of ten re-examined failures completes the attribution:
three never consulted the narration where it held the answer, three are long-video
localization failures (all on videos longer than 200\,s), and four are
comprehension failures at the model-capability ceiling. Of these classes only
long-video localization is mechanically addressable, and it is the most expensive;
we document it as the principal residual and leave it outside the scope of this
work.

\section{Transfer Check: AGQA}
\label{sec:agqa}

To verify the design is not MMBench-Video-specific, we ran the agent unchanged on a
70-question English subset of AGQA~\cite{grundemclaughlin2021agqa} (10 videos
$\times$ 7 questions; single run). AGQA's protocol differs from the main benchmark:
answers are scored both by exact match (EM) against the gold string and by an LLM
judge (\texttt{qwen-plus-latest}, 0/0.5/1 against key facts).

\begin{table}[t]
  \centering
  \small
  \caption{AGQA transfer results (70 questions, single run; LLM-judge / exact-match
  accuracies).}
  \label{tab:agqa}
  \begin{tabular}{@{}lccccc@{}}
    \toprule
    Metric & Overall & binary & open & sequencing & duration \\
    \midrule
    LLM-judge & 0.436 & 0.625 & 0.190 & 0.286 & \textbf{0.682} \\
    Exact match & 0.386 & 0.625 & 0.190 & 0.357 & 0.273 \\
    \bottomrule
  \end{tabular}
\end{table}

The profile transfers rather than the headline number
(\cref{tab:agqa}): the agent is strongest exactly where its design aims ---
\emph{duration} questions (0.682 under the judge), which require localising and
comparing event extents --- and weakest on open-vocabulary naming (0.190) and
sequencing (0.286), both dominated by the backbone's generation habits rather than
by evidence access. The divergence between the two metrics on duration (0.682 vs.\
0.273) is itself a methodological finding: exact match systematically
mis-scores a generative agent's free-form answers (we observed both false
negatives and a confirmed false positive where EM credited a wrong option-matching
answer the judge correctly rejected), so we treat the LLM-judge column as primary
and report EM for transparency.

\section{Annotation Audit and Threats to Validity}
\label{sec:threats}

\textbf{Annotation audit.} To check whether the remaining headroom is real, we
audited 30 randomly sampled questions (seed 11), asking whether the gold answer is
supported by the video's evidence (summary plus full transcript): 29 of 30 (97\%)
are fully supported, 1 partially, 0 unsupported --- a normalised score ceiling of
$\approx$0.98. Against the agent's $1.984/3 \approx 0.66$, the distance to the
ceiling reflects genuine capability gaps (fine-grained visual comprehension,
long-video localization), not label noise. \emph{Proxy disclosure}: the audit was
performed by an LLM reading the summary and transcript; it is biased lenient, and
questions whose evidence is purely fine-grained visual detail may be over-credited
as supported.

\textbf{Threats to validity}, each with its mitigation:
\begin{enumerate}
  \item \emph{Subset size and composition.} 150 questions with TR/HL oversampling
  is not the official distribution; re-weighting by the official distribution
  moves totals by $\leq 0.04$ (\cref{sec:setup}), and no leaderboard comparison is
  claimed.
  \item \emph{Judge self-preference.} The as-run judge (qwen-max) shares a vendor
  with the evaluated models. The full gpt-4-turbo re-judge refutes this concern:
  every method's total moves by less than $0.015$, the agent's lead widens
  slightly ($0.257 \to 0.265$), and per-question agreement (0.76--0.81) defines
  the judge-noise band against which we read all sub-$\pm 0.1$ differences.
  \item \emph{Single benchmark.} All main conclusions rest on MMBench-Video; the
  AGQA check (\cref{sec:agqa}) probes transfer but is small and single-run.
  Reproduction on a second long-video benchmark is the principal item of future
  work.
  \item \emph{Hosted models behind rolling aliases.} As stated in
  \cref{sec:impl}, backbone and VLM cannot be pinned to snapshots; all runs
  were executed within a single access window (June--July 2026).
  \item \emph{Moderation-driven exclusions.} The frame-scaling comparison removes
  three deterministically moderation-blocked questions from all conditions
  (common-147, \cref{subsec:framescaling}); the full-150 accounting in
  \cref{app:tables} shows the exclusion is conservative toward the baselines.
\end{enumerate}

\section{Case Study}
\label{sec:case}

Two traces from the benchmark runs make the mechanisms of \cref{ch:method}
concrete. (Question IDs refer to the released raw results; texts are condensed
from the recorded answers.)

\textbf{Case A --- grounded ``no'' (mmbv-1185, HL, score 3).} Asked \emph{``Did
NIO rank the 20th place in the chart?''} about a market-capitalisation ranking
video, the agent's free \textsc{Search} over $L_0$ finds a company list in the
global summary --- without NIO. Under grounding rule~1 of \cref{subsec:grounding}
(absence in memory is not evidence of absence in the video), it does not answer
``no'' yet: it explores 120--150\,s, then 60--120\,s (3 tool calls, 12 frames
total), finding XPeng, Subaru, Rivian and others, and a caption noting that each
chart block carries an explicit rank number. It then reasons about where rank \#20
must appear given Tesla's position at the end, concludes the explored windows cover
that range, and only then answers: \emph{``No. NIO does not appear in the video's
ranking --- it is absent from the global summary's named list and from two
explored segments (60--120\,s and 120--150\,s) covering the range where rank \#20
would appear.''} The answer is a negative claim earned by targeted perception ---
precisely the behaviour the HL dimension rewards (\cref{subsec:dimensions}). In
one of the three runs a different exploration path yielded only a partial answer
(score 1), a concrete instance of the exploration-path variance behind the agent's
larger standard deviation (\cref{sec:main-results}).

\textbf{Case B --- localization failure (mmbv-1552, TR, score 0 in all runs).} On
a 308-second physics-education video, the question asks what happens \emph{next}
after a yellow ball flies toward an apple. The agent explores 180--240\,s, finds
the segment where the ball illustrates motion without friction, but never locates
the actual collision moment; it falls back to inference from the global summary
and constructs a plausible but wrong sequence. The failure is not fabrication ---
the answer explicitly marks its inference --- but \emph{mislocalization under a
bounded budget} on a long video: the class identified in \cref{sec:negative} as
the principal residual gap, beyond the reach of the cheap levers tested there.

Together the cases bound the contribution: confidence-driven, provenance-grounded
exploration converts free-retrieval misses into evidence-backed answers
(Case A), and the dominant remaining failure is finding the right window in long
videos, not reasoning about it once found (Case B).
